\chapter*{Előszó}
\addcontentsline{toc}{chapter}{Előszó}

Ez a könyv webalkalmazás-fejlesztőknek szól, akik most találkoznak először az \rwlang{}-gal, és akár a modellközpontú alkalmazástervezéssel is. Nem feltételezi, hogy az olvasó ORM-ek, domain modellek vagy capability-alapú rendszerek fogalmaiban gondolkodik. Az első fejezet ezért nem szintaxissal kezd, hanem azzal, hogy \emph{miből áll egy alkalmazás}: milyen üzleti adatokat kezel, milyen műveletek végezhetők rajtuk, ki végezheti ezeket a műveleteket, és hol lép be mindez a HTTP világába.

A könyv ajánlott olvasási sorrendje tudatosan fokozatos. Először felépítjük a model- és capability-szemléletet, majd lefordítjuk és telepítjük az RWLang programokat. Ezután egy szándékosan kicsi, A--Z alkalmazáson végigmegyünk a modeltől a formon és adatbázison át a HTTP válaszig. Csak ezután bontjuk szét részletesen a nyelvi, adatbázis-, auth-, fájl- és böngészős témákat.

Amikor az olvasó már normál alkalmazást tud írni, következik a nagyobb projektek szervezése, egy wiki/CMS példa, a külső integrációk, a tesztelés, majd egy összefoglaló security-hardening fejezet. Ez a fejezet szándékosan a production üzemeltetés elé került: előbb legyen világos, mit garantál már a nyelv és a platform, és csak utána döntsünk skálázásról, megfigyelésről, deploymentről és recoveryről. A könyv végén a CLI és a teljes \code{server.toml} referencia már nem tananyagként, hanem visszakereshető operátori segédletként szolgál.

\section*{Hogyan olvasd a könyvet?}

Első olvasáskor ne próbáld megjegyezni az összes konfigurációs kulcsot vagy production limitet. Az első cél az, hogy három kérdésre mindig tudj válaszolni:

\begin{enumerate}
  \item \emph{Milyen typed adatot kezel ez a feature?}
  \item \emph{Milyen művelet olvassa vagy módosítja ezt az adatot?}
  \item \emph{Milyen authority vagy capability engedi ezt a műveletet?}
\end{enumerate}

Ha ez a három kérdés tiszta, a route, a form, az SQL, az auth és a konfiguráció már ugyanannak a modellnek a különböző határai lesznek, nem egymástól független trükkök.

A könyvben végig hasznos szétválasztani, mit bizonyít az RWLang, és mit kell továbbra is az alkalmazás tervezőjének eldöntenie:

\begin{center}
\begin{tabularx}{\textwidth}{@{}Y Y@{}}
\toprule
\textbf{RWLang/platform felelőssége} & \textbf{Alkalmazás/operator felelőssége} \\
\midrule
Typed boundaryk, escaping, capability check, bounded erőforrások, secure session mechanika & Üzleti szabályok, object ownership, tenant-membership policy, retention és recovery célok \\
Fail-closed compiler/runtime invariantok & A helyes route policy, permission, critical-operation contract és production topológia kiválasztása \\
Biztonságos SQL-, redirect-, upload-, egress- és crypto-primitívek & Annak eldöntése, mikor megengedett az üzleti művelet és mi történjen domain conflict esetén \\
\bottomrule
\end{tabularx}
\end{center}

A secure-by-construction modell egész plumbing-osztályokat vesz le a fejlesztőről; az alkalmazás üzleti authorization szabályait viszont nem találja ki helyette.

\subsection*{Ha sietsz}

A könyvet útvonalként használd, ne checklistként:

\begin{itemize}
  \item \textbf{első működő alkalmazás:} 1--5. fejezet;
  \item \textbf{autentikált CRUD vagy API:} folytasd az adatbázis-, auth-, fájl- és JSON-fejezetekkel, majd release előtt olvasd el a security-hardening fejezetet;
  \item \textbf{production üzemeltetés:} ehhez jön a cache/skálázás, observability, deployment, recovery és az operátori referencia.
\end{itemize}

\section*{Öt visszatérő szó}

A könyv néhány angol eredetű technikai szót következetesen használ. Már az elején érdemes ugyanazt érteni alattuk:

\begin{description}
  \item[\code{model}] fordító által ismert, névvel ellátott typed adatalak; nem ORM-osztály és nem automatikus adatbázisséma;
  \item[typed] az adat alakja és típusa a fordító számára ismert, nem csak futás közben derül ki;
  \item[capability] explicit hozzáférési lehetőség egy erőforráshoz vagy érzékeny művelethez, például \code{Db}, \code{Transaction} vagy \code{AppFs};
  \item[policy] a külső határon érvényes szabály, például auth, rate limit, CORS vagy trusted proxy;
  \item[runtime] a lefordított RWLang programot végrehajtó szerveroldali környezet, amely a limiteket és capability-határokat is kikényszeríti.
\end{description}

A későbbi fejezetek ezeket pontosítják, de a fenti jelentés első olvasásra elég.

\section*{A könyv konvenciói}

\begin{itemize}
  \item A \code{main.rw} az alkalmazás belépési pontja.
  \item A production konfiguráció elsődleges felülete a trusted TOML fájl.
  \item Secret érték nem kerül parancssorba vagy plaintext TOML mezőbe; a konfiguráció \code{*_file} hivatkozásokat használ.
  \item SQL szöveg compiler-owned, felhasználói érték csak typed bind paraméterként kerülhet lekérdezésbe.
  \item HTML-ben a \code{{ ... }} interpoláció escape-elt; raw HTML escape hatch nincs.
  \item Üzleti módosítás transactionben történik.
\end{itemize}

A könyv a V1 surface-t célozza. Ha egy capability nincs benne ebben a surface-ben, azt nem kerülőúttal kell megvalósítani, hanem külön nyelvi/runtime feature-ként.
